\section{An independent joint eigenvalue tuple}\label{sec:spectral}

We now combine the descendant lower bound with the recurrence to show
that the quantum divisors vary in all $r$ numerical directions.
Let
\[
 L=\overline{\mathbb C(q_1,\ldots,q_r)}.
\]
The commuting matrices $A_1,\ldots,A_r$ can be simultaneously put in
upper triangular form over $L$. A \emph{joint diagonal tuple} is the
$r$-tuple of diagonal entries at one position in such a form.

\begin{proposition}\label{prop:spectral}
At least one joint diagonal tuple $(\lambda_1,\ldots,\lambda_r)$ of the
quantum divisor matrices has coordinates algebraically independent over
$\mathbb C$.
\end{proposition}

We first normalize the descendant recurrence to obtain a uniform
exponential upper bound. We then show that a polynomial relation among
every joint tuple would force superexponential decay along a free degree,
contradicting Proposition~\ref{prop:descendant-bound}.

\subsection{Factorial normalization}

For $d_\beta\ge0$ and $0\le k\le n$, put
\[
 w_{\beta,k}=(d_\beta+k)!\,v_{\beta,k},
 \qquad w_\beta=\sum_{k=0}^n w_{\beta,k},
\]
where $v_{\beta,k}$ is the $H_k$ component. Put $w_\beta=0$ when
$d_\beta<0$. For a sequence $u$ indexed by $\mathbb Z^r$, let
$(E^\gamma u)_\beta=u_{\beta-\gamma}$ and set
$A_j(E)=\sum_\gamma A_{j,\gamma}E^\gamma$.
These are finite sums of shifts, and they commute.

If $A_{j,\gamma}$ sends an input component $H_l$ to $H_k$, then
$l=k-1+d_\gamma$. Thus
\[
 d_{\beta-\gamma}+l=d_\beta+k-1.
\]
Multiplying the $H_k$ component of \eqref{eq:recurrence} by
$(d_\beta+k-1)!$ gives, for $d_\beta>0$,
\begin{equation}\label{eq:normalized-recurrence}
 (A_j(E)w)_{\beta,k}
       =\frac{\beta_j}{d_\beta+k}\,w_{\beta,k}.
\end{equation}
Every nonzero input term has a nonnegative factorial argument. We do not
use this formula at degree zero.

Fix a norm on $H$. There is a constant $B>1$ such that
\begin{equation}\label{eq:exponential-bound}
 \|w_\beta\|\le B^{d_\beta+1}\qquad(d_\beta\ge0).
\end{equation}
To prove this, let $C_*=(\sum_jc_jD_j)\cup(-)$, and let $T_d$ multiply
$H_k$ by $1+k/d$ for positive integers $d$. Taking the linear
combination of \eqref{eq:normalized-recurrence} with coefficients $c_j$
yields
\[
 w_\beta=(I-T_{d_\beta}C_*)^{-1}T_{d_\beta}
              \sum_{\gamma>0}d_\gamma S_\gamma w_{\beta-\gamma}.
\]
Since $T_dC_*$ raises grading, the inverse $(I-T_dC_*)^{-1}$ is the
finite sum $\sum_{a=0}^n(T_dC_*)^a$. Its norm, and that of $T_d$, are uniformly
bounded for $d\ge1$. The sum uses finitely many positive shifts, each of
which strictly decreases $d$. Induction on the nonnegative integer $d$
proves \eqref{eq:exponential-bound}, taking $B$ large enough for the
finite set of coefficients and the initial value $w_0=n!\pt$.
All other degree-zero coefficients are zero.

\subsection{A relation would force too much decay}

\begin{proof}[Proof of Proposition~\ref{prop:spectral}]
Suppose, contrary to the proposition, that every joint diagonal tuple
is algebraically dependent over $\mathbb C$. There are finitely many
such tuples. Multiplying a nonzero polynomial relation for each of them
gives a nonzero $P\in\mathbb C[T_1,\ldots,T_r]$ vanishing on all of them.
If $m=\dim H$, simultaneous triangularization gives
$P(A_1,\ldots,A_r)^m=0$.
Substitute $(t^{c_j}q_j)_j$ and use \eqref{eq:homogeneity}. It follows
that $P(tA_1,\ldots,tA_r)^m=0$. The leading coefficient in $t$ gives
\begin{equation}\label{eq:homogeneous-relation}
 Q(A_1,\ldots,A_r)=0
\end{equation}
for a nonzero homogeneous polynomial $Q$: the $m$th power of the top
homogeneous part of $P$.

By Lemma~\ref{lem:free-dense}, choose a nonconstant free degree $\eta$
with $Q(\eta)\ne0$. Set $p_* =\eta/d_\eta$. Homogeneity gives
$Q(p_*)\ne0$. Choose a small bounded neighborhood of $p_*$ on whose
closure $|Q|$ is bounded below by a positive constant.
The substitution $q^\gamma\mapsto E^\gamma$ is a homomorphism from
Laurent polynomials to constant-coefficient shift operators. Applying it
entrywise to \eqref{eq:homogeneous-relation} gives
$Q(A_1(E),\ldots,A_r(E))=0$. Since these operators commute,
polynomial differencing gives
\begin{equation}\label{eq:polynomial-difference}
 Q(p)I=\sum_{j=1}^rB_j(p,E)(p_jI-A_j(E)).
\end{equation}
Here $p=(p_1,\ldots,p_r)$ consists of scalar parameters and every $B_j$
is a matrix polynomial in $p$ and finitely many shifts. This follows by
successively replacing the commuting variables in $Q(p)-Q(A(E))$.

Apply \eqref{eq:polynomial-difference} to $w$ at an index $\beta$,
with $p=\beta/d_\beta$ held fixed throughout this application.
For any shift $\delta$ that occurs and any $k$, the relevant component
of the rightmost factor is
\[
 \begin{aligned}
 ((p_jI-A_j(E))w)_{\beta-\delta,k}
 &=\left(\frac{\beta_j}{d_\beta}
    -\frac{\beta_j-\delta_j}{d_\beta-d_\delta+k}\right)
       w_{\beta-\delta,k}\\
 &=\frac{p_j(k-d_\delta)+\delta_j}{d_\beta-d_\delta+k}
       w_{\beta-\delta,k}.
 \end{aligned}
\]
When $p$ stays in the chosen neighborhood and $d_\beta$ is large,
the last scalar factor is uniformly $O(d_\beta^{-1})$: its numerator
is bounded, and its denominator differs from $d_\beta$ by a bounded amount.
The finitely many coefficient matrices in $B_j$ are uniformly bounded
there. Hence there are $C>0$ and a finite set of shifts $U$ such that
\begin{equation}\label{eq:small-shift-bound}
 \|w_\beta\|\le\frac C{d_\beta}
                     \max_{\delta\in U}\|w_{\beta-\delta}\|.
\end{equation}
If the right side of \eqref{eq:polynomial-difference} is the zero
operator, it directly implies $w_\beta=0$ in this neighborhood, which
already contradicts the descendant lower bound. We may therefore take
$U$ nonempty, including the zero shift if it occurs.

Take $\beta=s\eta$ for positive integers $s$. For sufficiently small
fixed $\epsilon>0$, every index reached by at most
$\lfloor\epsilon d_\beta\rfloor$ successive shifts from $U$ has degree
in $[d_\beta/2,2d_\beta]$ and direction in the chosen neighborhood.
Indeed each step changes both the vector and its degree by bounded
amounts; the total change is at most a constant times
$\epsilon d_\beta$. Choose $\epsilon$ first to keep the degree above
$d_\beta/2$, then to keep the normalized direction close to $p_*$.
For large $s$ all these indices also exceed the fixed lower threshold
for \eqref{eq:small-shift-bound}. Iteration, followed by
\eqref{eq:exponential-bound}, gives
\begin{equation}\label{eq:superexponential-bound}
 \|w_\beta\|\le
 \left(\frac{2C}{d_\beta}\right)^{\lfloor\epsilon d_\beta\rfloor}
 B^{2d_\beta+1}.
\end{equation}
This finite iteration does not require the shifts to decrease the
degree strictly, and it remains valid at indices with zero-extended
coefficients.

On the other hand, every positive multiple $s\eta$ is free: compose a
free map of degree $\eta$ with a degree-$s$ cover of $\mathbb P^1$.
By \eqref{eq:string-pairing} and
Proposition~\ref{prop:descendant-bound},
\[
 (w_\beta,1)
 \ge\frac{(d_\beta+n)!}{(a d_\beta)^{d_\beta-1}}
 \ge a\,d_\beta^{n+1}(ae)^{-d_\beta},
\]
where $a$ is fixed by the anticanonical embedding. The last inequality
uses $d!\ge(d/e)^d$. Pairing with $1$ is a bounded linear functional on
$H$, so this bounds $\|w_\beta\|$ below up to a fixed factor.
Its logarithm is bounded below by $-O(d_\beta)$, whereas the logarithm
of \eqref{eq:superexponential-bound} is
$-\epsilon d_\beta\log d_\beta+O(d_\beta)$. This contradiction proves
Proposition~\ref{prop:spectral}.
\end{proof}
